\begin{table}[H]
\caption{Total memory, $M_{\mathrm{total}}(D,N)=S\,D+T(N)$, from the measured
retained means of Table~\ref{tab:qcomem-validation60} and the measured transient
fits of Table~\ref{tab:transient} for $D$ warm documents and $N$ concurrent
requests.}
\label{tab:composition}
\centering\scriptsize
\setlength{\tabcolsep}{4.0pt}
\begin{tabular}{@{}rrrrr@{}}
\toprule
\multicolumn{5}{@{}l}{Warm-set crossover $D$: $5.13$ at $N=1$, $4.69$ at $N=2$, $3.79$ at $N=4$, $2.01$ at $N=8$}\\
Warm documents $D$ & Requests $N$ & Full-prefix (MiB) & \qcomem{} (MiB) & Ratio \\
\midrule
4 & 8 & 1978.6 & 1726.4 & $1.15\times$ \\
20 & 4 & 3493.3 & 1441.9 & $2.42\times$ \\
100 & 1 & 13893.2 & 1885.5 & $7.37\times$ \\
100 & 8 & 15057.2 & 2653.9 & $5.67\times$ \\
\bottomrule
\end{tabular}
\vspace{0.5mm}
\parbox{0.97\textwidth}{\scriptsize Every cell carries the three unmeasured
assumptions of Section~\ref{sec:composition}: $S$ linear in $D$ (no run held
more than one document resident); $T$ independent of $D$ (unverified); fits
valid only for $N\in[1,8]$, so no cell licenses a statement about $N>8$.  The
transient term is an allocator peak on this stack, is not additive with Store,
and fixes no admitted capacity.  No workload of the many-warm-document shape
these rows describe was executed.}
\end{table}
